%%%%%%%%%%%%
%
% $Autor: Wings $
% $Datum: 2019-03-05 08:03:15Z $
% $Pfad: Maintenance.tex $
% $Version: 4250 $
% !TeX spellcheck = en_GB/de_DE
% !TeX encoding = utf8
% !TeX root = manual 
% !TeX TXS-program:bibliography = txs:///biber
%
%%%%%%%%%%%%

\chapter{Wartung}
	\textbf{Allgemeiner Hinweis:}\\
	Für jegliche Wartungsarbeiten muss die Stromzufuhr unterbrochen werden. Die Stromzufuhr darf erst nach Beendigung der Wartungsarbeiten eingeschaltet werden \\ \\
	\textbf{Wartungshinweise:}
	\begin{enumerate}
		\item Nach längeren Stillstand des Wasserkreislaufs sollte das Wasser im Eimer ausgewechselt werden. Um das Wasser auszutauschen, wird der Deckel mit Acryl-Rohr abgenommen. Danach wird der Eimer entleert und nach Bedarf gereinigt. Dann kann man den Blumentopf mit frischem Wasser bis zur max. Füllhöhe befüllen.
		\item Um die Acrylrohre zu reinigen wird zuerst der Deckel abgenommen. Um das Rohr für den Schwimmer zu reinigen, wird das Acrylrohr vom Boden des Eimers abgeschraubt und mit dem Schwimmer aus dem Eimer genommen. Nach der Reinigung wird das Acrylrohr wieder im Eimer angeschraubt und der Schwimmer in das Acrylrohr gelegt. 
		Um das Acrylrohr für die Beförderung des Wassers zu Reinigen wird nach dem Entfernen des Deckels das Rohr von der Pumpe Losgeschraubt und durch den Deckel nach oben hin ausgezogen, der Wasserhahn kann ebenfalls losgeschraubt werden.
		\item Um den Ultraschallsensor zu reinigen, wird der Deckel abgenommen das Acrylrohr für den Wassertransport ebenfalls entnommen. Danach kann der Boden vom Deckel entfernt werden. Der Ultraschallsensor darf nur mit einem Mikrofasertuch und sanften Bewegungen gereinigt werden der Kontakt mit ist zu unterlassen! 
		\item Wenn der Schwimmer ersetzt oder gereinigt werden soll kann der Deckel abgenommen werden das Acrylrohr für den Schwimmer vom Boden des Eimers soll geschraubt und entfernt werden. Der Schwimmer kann dann Einfach aus dem Eimer entnommen werden. Um den Schwimmer wieder in Position zu bringen wird das entsprechende Acrylrohr am Boden des Eimers angeschraubt und der Schwimmer von oben in das Acrylrohr eingelegt, danach kann der Deckel wieder aufgesetzt werden.
	\end{enumerate}
	 


